% RAG Assignment report
%
% HOW TO USE THIS FILE
% Every table, number and figure below is real output from rag_rerank.ipynb.
% Every DRAFT block is a starting point -- rewrite it in your own words.
% Delete the \draft macro definition and all \draft calls before submitting.
%
% Build:  pdflatex report.tex

\documentclass[11pt,a4paper]{article}
\usepackage[margin=1in]{geometry}
\usepackage{booktabs}
\usepackage{xcolor}
\usepackage{hyperref}
\usepackage{listings}

% --- remove this macro (and every use of it) once you've written your text ---
% DRAFT text is a starting point only -- rewrite before submitting.
\newcommand{\draft}[1]{\medskip\noindent\textcolor{gray}{#1}\medskip}

\title{Baseline Retrieval vs.\ Cross-Encoder Re-Ranking:\\
       A Two-Stage RAG Study on the Gloomhaven Rulebook}
\author{Sampath Ravikanti}
\date{October 2026}

\begin{document}
\maketitle

\section{Setup}

\begin{tabular}{ll}
\toprule
Corpus          & Gloomhaven rules reference, 74{,}916 characters \\
Sectioning      & 39 headed sections detected from the layout \\
Chunking        & 550 characters, section-aware; 167 chunks \\
Chunk sizes     & min 160 / median 490 / max 571 characters \\
Bi-encoder      & \texttt{all-MiniLM-L6-v2} (384-dim, local) \\
Cross-encoder   & \texttt{cross-encoder/ms-marco-MiniLM-L-6-v2} \\
Generator       & \texttt{gemini-3.8-flash} \\
Vector store    & normalised matrix, cosine via one NumPy matmul \\
\bottomrule
\end{tabular}

\medskip
Both pipelines share stage-1 retrieval and call the same generation function with
an identical prompt string, so the only variable between Answer 1 and Answer 2 is
which three chunks are selected.

\draft{A rulebook was chosen over a policy or legal document deliberately. The failure
modes this assignment asks for -- keyword traps, exceptions, negation -- all require a document
that defines many near-synonymous terms precisely and then qualifies them. Gloomhaven separates
\emph{exhausted} from \emph{eliminated} and \emph{lost} from \emph{discarded}, and states most
rules as a general case followed by a restriction located paragraphs away. A regulatory PDF has
similar structure but spreads it over far more text, meaning a larger corpus for the same number
of usable test cases.}

\section{Architecture Comparison}

\draft{The two models differ in when the query and the document meet.

A bi-encoder passes each text through the transformer separately and mean-pools the token states
into one fixed vector. Chunk vectors depend on nothing but the chunk, so they can be computed
once at index time; a query is embedded at run time and compared by cosine similarity. The entire
interaction between query and chunk is a single dot product between two vectors produced in
ignorance of each other. Whatever mean-pooling discarded is unavailable at scoring time, and
pooling is lossy in a specific way: it is order-insensitive and dominated by tokens carrying the
most semantic weight. Function words that invert meaning -- \emph{not}, \emph{cannot},
\emph{except} -- contribute little relative to content words, which is why ``a stunned figure may
use items'' and ``a stunned figure cannot use items'' land close together in the space.

A cross-encoder produces no comparable vectors at all. It concatenates query and chunk into one
sequence and runs full self-attention over the pair, so every query token attends to every chunk
token. The output is a single relevance score for that specific pairing. This lets the model
condition on exact token alignment: \emph{cannot} in the chunk can attend directly to \emph{may}
in the query and change the score, which is impossible once both have been pooled.

The cost of that expressiveness is that nothing can be precomputed. A bi-encoder index is built
once; a cross-encoder score exists only after the query arrives and must be recomputed for every
candidate. This asymmetry is the entire justification for the two-stage design, and Section 6
quantifies it.}

\section{Query Set and Method}

Ten queries were written against rules located by reading the document first, then
phrased deliberately in everyday synonyms rather than the rulebook's own vocabulary
(``strike back'' for \emph{retaliate}, ``shut entryway'' for \emph{closed door}).
A gold chunk was identified by hand for each query, and a query counts as answered
only when the gold chunk appears in the top 3.

\draft{Phrasing matters more than it first appears. If queries reuse the document's own
vocabulary, lexical overlap alone ranks the correct chunk first and both pipelines succeed,
measuring nothing. Writing ``strike back'' instead of \emph{retaliate} and ``shut entryway''
instead of \emph{closed door} removes that shortcut and forces retrieval to work semantically --
the condition under which the architectures actually differ. Queries were written after reading
the rulebook and locating a genuinely tricky rule, rather than by guessing and hoping some would
fail.}

\section{Qualitative and Rank-Shift Analysis}

\subsection{Outcome across the query set}

\begin{tabular}{llcl}
\toprule
Verdict & & bi $\rightarrow$ ce & Query \\
\midrule
Fixed by re-rank    && 7 $\rightarrow$ 1 & am I allowed to skip part of what my card tells me to do? \\
Fixed by re-rank    && 6 $\rightarrow$ 2 & may a stunned figure drink a healing potion on its turn? \\
Fixed by re-rank    && 7 $\rightarrow$ 3 & what occurs if a monster has nothing it is able to aim at? \\
Artifact (see 4.4)  && 4 $\rightarrow$ 3 & may I stockpile two of the same piece of equipment? \\
Both correct        && 1 $\rightarrow$ 1 & does mending someone's injuries also cure the venom? \\
Both correct        && 2 $\rightarrow$ 1 & can a beast walk through a shut entryway? \\
Broken by re-rank   && 2 $\rightarrow$ 6 & do I get a bigger hand of cards as I grow more powerful? \\
Broken by re-rank   && 2 $\rightarrow$ 4 & do sealed entrances open when someone steps onto them? \\
Stage-1 miss        && --- & if I strike back at someone who hit me, does that count as attacking? \\
Stage-1 miss        && --- & is there a ceiling on restoring a wounded hero's vitality? \\
\bottomrule
\end{tabular}

\medskip
Three genuine rescues, two regressions, two failures that re-ranking could not reach.

\subsection{Rank shifts on the rescued queries}

\begin{tabular}{llcccc}
\toprule
Query & Gold chunk & bi-rank & ce-rank & bi-score & ce-score \\
\midrule
skip part of an action      & 17 & 7 & 1 & 0.330 & $-3.475$ \\
stunned figure, potion      & 43 & 6 & 2 & 0.344 & $-1.680$ \\
monster with no target      & 71 & 7 & 3 & 0.476 & $-1.628$ \\
\bottomrule
\end{tabular}

\medskip
Note the bi-encoder scores: on the first query the gold chunk scored 0.330 against
0.396 for the top-ranked distractor -- a gap of 0.066 across seven rank positions.

\draft{The score distributions are the most informative part of this table. On the first
query the gold chunk scored 0.330 while the top-ranked distractor scored 0.396 -- 0.066
separating rank 1 from rank 7, with all ten candidates inside a band roughly 0.07 wide. When
every topically-related chunk receives nearly the same similarity, ordering within that band is
determined by differences that correspond to nothing meaningful about relevance.

This is a predictable consequence of mean-pooled cosine similarity. Pooling maps any passage
about taking a turn into approximately the same region of the space, so chunks sharing a topic
share a score, and the metric has no resolution left to distinguish the chunk that answers the
question from six that merely discuss it.

Cross-encoder scores on the same candidates span $-3.5$ to $-11.3$ -- several points rather than
hundredths. Scoring each pair jointly, the model is not committed to placing all turn-related
chunks together, and separates them decisively.}

\subsection{Answer 1 vs Answer 2}

On the stun query the baseline did not hallucinate -- it correctly reported that it
could not answer:

\begin{quote}
\textbf{Answer 1 (baseline):} ``the answer is not contained within the text\ldots\ the
context does not mention the `stunned' condition or rules regarding drinking potions.''
\end{quote}

\begin{quote}
\textbf{Answer 2 (re-ranked):} ``No, a stunned figure cannot drink a healing potion on
its turn. \emph{STUN: The figure cannot perform any abilities or use items except, if a
character, to perform a long rest.}'' (chunk 43)
\end{quote}

\draft{The baseline's failure mode deserves attention, because it is not the one the
assignment anticipates. It did not hallucinate. It correctly determined that its three chunks did
not address the stun condition and said so.

That behaviour is a product of the prompt, which instructs the model to state explicitly when
context is insufficient rather than guess. Under that constraint, retrieval failure degrades into
abstention instead of fabrication -- the safer outcome, since a user told ``this is not in the
context'' knows to look elsewhere, while a user given a confident wrong answer does not.

It would be wrong to conclude the prompt solves the problem. The same three chunks under a less
careful prompt, or with a model more disposed to fill gaps, would plausibly have produced an
answer assembled from the general action rules -- fluent, specific and wrong. And abstention is
still failure: the user asked an answerable question about a rule present in the document and
received nothing. The prompt converted a wrong answer into a missing one. Retrieval quality sets
a ceiling on answer quality that prompt engineering can protect but cannot raise.}

\draft{The two regressions sharpen this. Both involve a general rule and a specific
instance of it. For the hand-size query the cross-encoder promoted chunks about levelling up and
card availability over the chunk stating that hand size is fixed at the value printed on the
character mat; for the pressure-plate query it preferred passages about doors and room reveals
over the specific rule on sealed entrances. In both cases the demoted chunk answers the question
tersely, in a single clause, while the promoted chunks discuss the topic at greater length. A
re-ranker trained on web passage relevance appears to reward elaboration on the query's subject,
which is usually a sound heuristic and is wrong exactly when the answer is a one-line exception.
The failures are not random: they are the same mechanism that produces the rescues, applied where
it does not fit.}

\subsection{Failure categories covered}

\begin{tabular}{p{3.2cm}p{7.2cm}p{3.4cm}}
\toprule
Spec category & Query & Evidence \\
\midrule
Keyword trap / distractor &
am I allowed to skip part of what my card tells me to do? &
Baseline top 3 (chunks 18, 11, 16) are saturated with \emph{card}, \emph{action},
\emph{turn}; the gold chunk 17 phrases the rule as ``free to not perform any part''.
Rank 7 $\rightarrow$ 1. \\
\addlinespace
Exceptions \& corner cases &
\textcolor{red}{[not yet isolated -- see note below]} & --- \\
\addlinespace
Negation \& constraints &
may a stunned figure drink a healing potion on its turn? &
Gold rule is stated as \emph{cannot\ldots except}; baseline abstains, re-rank
quotes it. Rank 6 $\rightarrow$ 2. \\
\addlinespace
Buried chunk (ranks 4--10 promoted) &
what occurs if a monster has nothing it is able to aim at? &
Gold chunk 71 at rank 7, promoted to 3. Also observed at ranks 6, 7 and 4 on the
other queries. \\
\bottomrule
\end{tabular}

\medskip
\textcolor{red}{\textbf{TODO:} the exceptions/corner-case category currently shares the
stun query with negation. Add a dedicated query where the baseline returns a general
rule and the governing ``except when\ldots'' sits in a different chunk, then fill the
row above.}

\subsection{A measurement artifact worth reporting}

The duplicate-items query appears in the table as fixed, but it is not. The rulebook
states the rule twice in different words:

\begin{quote}
``You cannot own more than 1 copy of any item card.'' \\
``No character may own duplicates of the same item.''
\end{quote}

The gold substring matched only the first. The baseline's top 3 contained the second,
so its answer was correct while being scored wrong.

\draft{Any evaluation testing for a fixed substring inherits this problem, and the
bias is not symmetric. A document stating a rule once in the gold string's words is scored
correctly; a document that paraphrases itself produces false failures for whichever pipeline
retrieved the paraphrase. Since the measurement was built around cases where the baseline was
expected to fail, errors of this kind flatter the re-ranker. Two fixes exist: accept a set of
acceptable phrasings per query, or abandon chunk identity as the criterion and judge the
generated answer against the rule itself. The second is more work and more faithful, since
retrieving a different chunk supporting the same correct answer is not a retrieval failure in any
sense a user would recognise.}

\subsection{Sensitivity to chunk size}

\begin{tabular}{ccccc}
\toprule
Chunk size & Chunks & Fixed & Broken & Stage-1 misses \\
\midrule
450 & 219 & 2 & 0 & 3 \\
550 & 167 & 4 & 2 & 2 \\
700 & 126 & 1 & 1 & 2 \\
\bottomrule
\end{tabular}

\draft{This table is the most important caveat in the report, and it bounds the headline
result substantially.

At 700 characters the measured benefit nearly disappears: one rescue instead of four. The reason
is mechanical. Larger chunks are more likely to contain both a general rule and the exception
qualifying it, so the baseline retrieves a passage that happens to include the answer and stops
failing -- there is less for the re-ranker to repair. At 450 characters the opposite happens:
stage-1 misses rise to three, because rules fragment across chunk boundaries and no single chunk
holds enough context to score highly.

The consequence is that the benefit reported here is partly a function of a chunking parameter
chosen by the author. A different reasonable choice would have produced a visibly different
conclusion. Re-ranking is most valuable in a particular regime -- chunks small enough to isolate
the answer, and a corpus where many chunks are topically similar. Both hold at 550 characters on
this document; neither is universal. A report presenting only the 550 row would be overstating a
general claim on the strength of a tuned one.}

\section{Latency Profiling}

\begin{tabular}{lccc}
\toprule
& Baseline & With re-ranking & Overhead \\
\midrule
Mean over 10 queries & 8.2 ms & 169.8 ms & $20.7\times$ \\
Range                & 6.8--10.2 ms & 144.2--216.5 ms & 15.6--27.1$\times$ \\
\bottomrule
\end{tabular}

\medskip
Measured on CPU. Stage-1 cost is one query embedding (7.8\,ms) plus a
$167 \times 384$ matmul. The cross-encoder costs 14.6\,ms per (query, chunk) pair.

\draft{An overhead of $20.7\times$ sounds severe and is almost irrelevant in context.
The absolute figures are 8.2\,ms and 169.8\,ms, and both feed an LLM generation step taking one
to three seconds. Re-ranking adds roughly 160\,ms to a pipeline whose latency is dominated by a
cost an order of magnitude larger. The multiplier alarms only when read without units: it
compares a very fast operation to a slightly less fast one. Where re-ranking does become
expensive is throughput under concurrent load, since the cross-encoder occupies the CPU for
160\,ms of real compute while the matmul does not.}

\section{Scaling: Why Not Cross-Encode Everything?}

At a measured 14.6\,ms per pair:

\begin{tabular}{lr}
\toprule
Scope & Time per query \\
\midrule
Top-10 re-rank (what we do) & 0.15 s \\
This corpus, all 167 chunks & 2.44 s \\
A 100k-chunk corpus         & 1{,}458.66 s (24 min) \\
A 1M-chunk corpus           & 14{,}586.64 s (4.1 hours) \\
\bottomrule
\end{tabular}

\draft{The answer follows from where the query enters the computation.

For the bi-encoder, chunk vectors depend only on chunks. They are computed once when the index is
built and reused unchanged by every subsequent query, so per-query work is one embedding plus a
similarity search -- which an ANN index makes sublinear in corpus size. Adding a million chunks
increases build time and memory but barely changes query latency.

For the cross-encoder the score is a function of the pair. Nothing about a chunk's relevance can
be known before the query exists, so no work amortises across queries and cost is strictly linear
in candidates scored. The measured 14.6\,ms per pair makes this concrete: scoring this 167-chunk
corpus exhaustively takes 2.44 seconds per query, and a 100k-chunk corpus would take 24 minutes.
Not merely slow -- the wrong shape of computation to apply at corpus scale.

Two-stage retrieval resolves the asymmetry by letting each model do what it is cheap at. The
bi-encoder reduces the corpus to 10 candidates at index-time cost; the cross-encoder adjudicates
those 10 in 0.15 seconds, approximating cross-encoder precision at close to bi-encoder cost.

The arrangement has a hard limit, which this data shows plainly. The re-ranker can only reorder
what stage 1 retrieved. Two queries failed in both pipelines because the gold chunk never entered
the top 10 -- the retaliate and heal-cap queries, both vocabulary mismatches sharing almost no
content words with the governing rule. No improvement to the re-ranker could have recovered them.
Stage-1 recall is the ceiling; stage 2 only determines how much of it is realised.}

\section{Conclusion}

\draft{The experiment compared two pipelines differing in exactly one step, holding
corpus, chunking, prompt and generator constant. Across ten queries, cross-encoder re-ranking
produced three genuine rescues, two regressions, and left two failures untouched because the gold
chunk never reached the candidate set. One further apparent rescue proved an artifact of
substring scoring.

Three of ten is a real but modest effect, and two qualifications bound it. The chunk-size
ablation shows the benefit is largely specific to a size that isolates rules from their
exceptions; at 700 characters it nearly vanishes. And the regressions are not noise -- they stem
from the same preference for elaboration that produces the rescues, which makes re-ranking a
trade rather than a free improvement, its net value depending on how often a corpus states
answers tersely.

What the results show most clearly is the compression of bi-encoder scores. Ten candidates within
a 0.07 band, the correct chunk at rank 7 and 0.066 below a distractor, is not a near miss but an
indication that cosine similarity over pooled embeddings had no resolution left to make the
distinction. The cross-encoder separated the same candidates across several points. This is the
mechanism behind every rescue observed, and the concrete form of the architectural difference
described in Section 2.

Given more time the clearest improvement would target stage 1 rather than stage 2. Both
unreachable failures were vocabulary mismatches, precisely what a lexical signal handles well; a
hybrid BM25 and dense stage would likely have placed both gold chunks in the top 10 and handed
them to a re-ranker already shown capable of promoting them. Raising recall at stage 1 lifts the
ceiling, whereas further work on stage 2 only redistributes what lies beneath it.}

\appendix
\section{Reproducing}

\begin{lstlisting}[basicstyle=\small\ttfamily]
curl -sL https://www.orderofgamers.com/downloads/Gloomhaven_v1.1.pdf \
     -o gloomhaven_rulebook.pdf
export GEMINI_API_KEY="..."
jupyter notebook rag_rerank.ipynb
\end{lstlisting}

Code: \url{https://github.com/Sampath7099/RAG-assignment}

\end{document}
